\chapter{Desenvolvimento da Ferramenta}
\label{cap:desenvolvimento}

% -----------------------------------------------------------------
% Objetivo do capítulo:
% Mostrar a implementação concreta do que foi descrito no
% Capítulo 3. É aqui que entram diagramas de arquitetura,
% pseudocódigo dos indicadores, exemplos de saída, decisões de
% engenharia e ajustes feitos durante o desenvolvimento.
%
% Referência interna:
%   docs/pipeline/README.md
%   docs/pipeline/resumo_algoritmos.md
%   docs/pipeline/resumo_metodo_estatistico_global.txt
% -----------------------------------------------------------------


\section{Visão geral da arquitetura}
\label{sec:arquitetura}

% Conteúdo a desenvolver:
%   - Nome interno da ferramenta (WireSense) e organização em
%     dois pilares (Pilar MP e Pilar EQ).
%   - Diagrama da arquitetura (referenciar imagem em
%     docs/diagramas/).
%   - Papel de cada pilar e como se comunicam.


\section{Pipeline de dados}
\label{sec:pipeline}

% Conteúdo a desenvolver:
%   - Sequência de etapas, das tabelas brutas ao diagnóstico.
%   - Mapa entre scripts implementados e etapas do pipeline.
%   - Pontos de integração com as bases existentes.


\section{Implementação do módulo de anomalia (Pilar EQ)}
\label{sec:impl_eq}

% Conteúdo a desenvolver:
%   - Conjunto de 14 sensores e justificativa técnica da
%     escolha.
%   - Treinamento e ajuste do Isolation Forest.
%   - Critérios de SEM_MODELO e fallback ao Pilar MP.
%   - Exemplos de saída do módulo.


\section{Implementação do módulo estatístico (Pilar MP)}
\label{sec:impl_mp}

% Conteúdo a desenvolver:
%   - Indicadores calculados (taxas, exposição, risco relativo,
%     cruzamento entre máquinas).
%   - Decisão metodológica registrada em 12/05: o Pilar MP
%     conta apenas rompimentos, com bolha removida do
%     numerador.
%   - Regra MODERADA: o Pilar MP é acionável, o ground truth
%     do relatório de devoluções é reservado para validação
%     externa.
%   - Exemplos de saída do módulo.


\section{Regras de integração e diagnóstico final}
\label{sec:regras_integracao}

% Conteúdo a desenvolver:
%   - Lógica de combinação entre Pilar EQ e Pilar MP.
%   - Categorias de origem provável.
%   - Forma de apresentação do diagnóstico (texto, indicadores,
%     evidência de apoio).


\section{Decisões de engenharia e ajustes iterativos}
\label{sec:decisoes_engenharia}

% Conteúdo a desenvolver:
%   - Principais decisões tomadas durante o desenvolvimento.
%   - Compromissos entre cobertura, precisão e
%     interpretabilidade.
%   - Pontos de extensibilidade da ferramenta.
